% !TEX root = ../thesis.tex

\chapter{Related Work}
\label{chp:related}

\section{Deep Reinforcement Learning for Legged Locomotion}

Learning locomotion controllers for legged robots via deep RL has become
the dominant paradigm.
Lee~\emph{et al.}~\cite{lee2020learning} demonstrated that a recurrent
neural network trained in simulation could learn robust quadrupedal walking
over stairs and slippery surfaces, exploiting privileged observations during
training that are unavailable at deployment.
Kumar~\emph{et al.}~\cite{kumar2021rma} introduced the Rapid Motor
Adaptation (RMA) framework, in which a teacher policy with access to
privileged environment parameters trains a student adaptation module that
infers latent terrain properties from a window of proprioceptive history.
The two-stage distillation approach---teacher $\rightarrow$ student---is
widely adopted and is inherited in RoboDuet for the arm policy.

Margolis and Agrawal~\cite{margolis2023walk} proposed ``Walk These Ways''
(WTW), a single gait-conditioned locomotion policy that covers a wide range
of speeds and gait styles through a parameterised command space.
WTW serves as the locomotion backbone in RoboDuet and, by extension, in this
project.  Zero-velocity locomotion commands are used in our demonstrations to
keep the quadruped stationary.

\section{Robotic Arm Grasping and Manipulation}

End-effector reaching and grasping have been studied extensively in
fixed-base settings.
Deep RL methods such as those of OpenAI~\cite{andrychowicz2020dexterous}
and continuous control benchmarks~\cite{tassa2018dmcontrol} have shown that
model-free policies can solve complex multi-joint manipulation tasks given
sufficient simulation data.
For goal-conditioned reaching, spherical ($(l, p, \psi)$) or Cartesian
command representations can be used; RoboDuet adopts a spherical
parameterisation that naturally encodes reach distance and elevation angle,
which aligns well with the kinematic range of the Z1 arm.

The Unitree Z1 arm features six revolute joints:
\texttt{z1\_waist}, \texttt{z1\_shoulder}, \texttt{z1\_elbow},
\texttt{z1\_wrist\_angle}, \texttt{z1\_forearm\_roll},
\texttt{z1\_wrist\_rotate}, plus a passive gripper.
Its workspace extends approximately 0.3--0.77\,m in reach and $\pm$81\,$^\circ$
in elevation, enabling it to place the end-effector at a height of
$0.38 + 0.77\sin(0.45\pi) \approx 0.93$\,m when fully extended upward.

\section{Whole-Body Loco-Manipulation}

Combining locomotion and manipulation on a single platform has attracted
growing research attention.
Fu~\emph{et al.}~\cite{fu2023deepwhole} addressed whole-body control for a
wheeled mobile manipulator, learning a unified policy via RL that coordinates
all body joints.
Cheng~\emph{et al.}~\cite{cheng2023legs} demonstrated that quadruped leg
movements can themselves be used as a form of manipulation, using contact
forces at the feet to push objects.
Zimmermann~\emph{et al.}~\cite{zimmermann2021go} mounted an arm on a Go1 and
used classical trajectory optimisation to plan arm motions while the legs
provided stability.
These works highlight a recurring design choice: \emph{centralised} control
(one policy outputs all joint commands) vs.\ \emph{decentralised} control
(separate locomotion and manipulation policies that interact implicitly).

\section{The RoboDuet Framework}

Pan~\emph{et al.}~\cite{pan2024roboduet} proposed RoboDuet, the decentralised
two-policy framework that forms the basis of this project.
RoboDuet trains two cooperative policies simultaneously: a dog policy
$\pi_\text{dog}$ controlling the 12 leg joints, and an arm policy
$\pi_\text{arm}$ controlling the 6 arm joints.
Neither policy has explicit access to the other's action output; they
interact only through the shared body dynamics.
Key features include:
\begin{itemize}
  \item \textbf{Spherical arm command space} $(l, p, \psi, \alpha, \beta,
        \gamma)$: reach length, pitch, yaw, and optional end-effector roll,
        pitch, yaw for full 6D pose tracking.
  \item \textbf{Adaptation module}: a learned mapping from proprioceptive
        observation history to a low-dimensional latent vector that
        approximates privileged information, enabling sim-to-real transfer
        without explicit environment identification.
  \item \textbf{Zero-shot cross-morphology transfer}: the authors report
        successful transfer from Go1+Z1 to a different quadruped (B2+Z1)
        without retraining, which motivates the present adaptation work.
  \item \textbf{Reward decomposition}: the dog reward includes locomotion
        tracking terms ($r_\text{lin}$, $r_\text{ang}$) and whole-body
        terms ($r_\text{orient}$, $r_\text{hip}$); the arm reward is
        dominated by the manipulation command tracking term.
\end{itemize}

The original paper demonstrates 23\% improvement in success rate over a
single-policy baseline on challenging loco-manipulation tasks.
Our work validates RoboDuet's transfer claim on the heavier B2Z1 platform,
develops a reproducible demonstration harness, and documents the engineering
decisions required for successful adaptation.
